\subsection{Neues Projekt} \label{subsec:Neues Projekt}
\newcommand{\CCBDarlehen}{75.000,00}
\newcommand{\CCBZinssatz}{3.5}
\newcommand{\CCBLaufzeit}{60}
\newcommand{\CCBAnnuitaet}{1.364,38}
\newcommand{\CCBGesamtzinsen}{6.862,85}
\newcommand{\CCBGesamttilgung}{75.000,00}


\paragraph{Grundidee} \label{para:Grundidee}

\par Die \gls{stz} plant, ein neues Projekt zu starten, um ihr Immobilienportfolio zu erweitern. Das Projekt umfasst die Akquisition einer neuen Wohneinheit in München, um von der positiven Entwicklung des Immobilienmarktes in der Region zu profitieren. Das Ziel ist dabei eine Wohneinheit in attraktiver Lage zu erwerben, die ein hohes Entwicklungspotenzial aufweist, um langfristige Wertsteigerungen zu erzielen. Das Projekt soll innerhalb eines festgelegten Zeitrahmens und Budgets durchgeführt werden, um eine effiziente Umsetzung sicherzustellen. 

\paragraph{Finanzierung}

\par Zur Finanzierung des Projekts wird die \gls{stz} eine Kombination aus Eigenkapital und Fremdkapital in Betracht ziehen. Um das Eigenkapital zu erhöhen, wird überlegt durch den Kommanditisten Sebastian Seitz eine Kapitalerhöhung oder ein Privatdarlehen durchzuführen, um zusätzliche Mittel für die Akquisition bereitzustellen. 

\begin{itemize}
    \item \textbf{Kapitalerhöhung}: Durch die Privateinlage des Kommanditisten Sebastian Seitz kann das Eigenkapital der \gls{stz} erhöht werden, was die finanzielle Stabilität verbessert und die Kreditwürdigkeit erhöht.
    \item \textbf{Privatdarlehen}: Alternativ könnte Sebastian Seitz der \gls{stz} ein Privatdarlehen gewähren, um die benötigten Mittel für die Akquisition bereitzustellen. Dies könnte eine flexible Finanzierungsoption sein, die es der \gls{stz} ermöglicht, die Rückzahlung an die finanziellen Möglichkeiten anzupassen.
    \item \textbf{Fremdprivatdarlehen}: Es könnte in Erwägung gezogen werden, ein Fremdprivatdarlehen von einer externen Kreditgeberin (Clara Craffonara) aufzunehmen, um die Finanzierung des Projekts zu unterstützen. Dies könnte eine zusätzliche Finanzierungsquelle sein, um die benötigten Mittel für die Akquisition bereitzustellen.
\end{itemize}

\par Ein Fremdprivatdarlehen von Clara Craffonara könnte eine attraktive Option sein, da es möglicherweise günstigere Konditionen bietet als traditionelle Bankkredite. Es ist jedoch wichtig, die Bedingungen des Darlehens sorgfältig zu prüfen und sicherzustellen, dass sie mit den finanziellen Zielen und Möglichkeiten der \gls{stz} übereinstimmen. Beispielhaft wird in \fullref{tab:tilgungsplan} ein Tilgungsplan für ein hypothetisches Fremdprivatdarlehen von Clara Craffonara aussehen könnte. Das Darlehen beträgt \CCBDarlehen{} EUR bei einem Zinssatz von \CCBZinssatz{}\,\% über \CCBLaufzeit{} Monate. Die monatliche Rate beläuft sich auf \CCBAnnuitaet{} EUR.
Die Gesamtzinslast beträgt \CCBGesamtzinsen{} EUR.


\newpage

\par Grundsätzlich sind aus meiner Sicht noch die Folgenden Punkte zu klären, um den Finanzrahmen des Projekts \glqq Immobilieneerwerb \grqq  zu finalisieren:

\begin{enumerate}
    \item Cashflow-Analyse aus Bestand erheben
    \item Optionserarbeitung für die Finanzierung des Projekts (Kapitalerhöhung, Privatdarlehen, Fremdprivatdarlehen)
    \item Kreditrahmen mit Bank klären 
    \item Konditionen für Privatdarlehen mit Sebastian Seitz klären
    \item Konditionen für Fremdprivatdarlehen mit Clara Craffonara klären
    \item Entscheidung über Finanzierungsstruktur treffen
\end{enumerate}


\par Als möglicher Rahmen zur Eigenmittelaufnahme könnten die folgenden Schritte dienen:

\begin{itemize}
    \item Mögliche Kapitalerhöhung durch den Kommanditisten Sebastian Seitz von $75.000$ EUR prüfen, um das Eigenkapital der \gls{stz} zu erhöhen.
    \item Möglichs Privatdarlehen von Sebastian Seitz in Höhe von $75.000$ EUR prüfen, um die benötigten Mittel für die Akquisition bereitzustellen.
    \item Mögliches Fremdprivatdarlehen von Clara Craffonara in Höhe von $75.000$ EUR prüfen, um die Finanzierung des Projekts zu unterstützen.
\end{itemize}